% GENERATED FILE. Edit canonical lessons, then run npm run generate:books.
% canonical-source-hash: fnv1a64:a0e12dc4ac39b9d1
% canonical-lessons: RU-C532-oboi, RU-C532-kleit, RU-C532-sverlit, RU-C532-drel, RU-C532-otvyortka, RU-R532-review-525-532

\chapter{Wallpaper and Drills}
\label{ch:ru-wallpaper-and-drills}
\hlchaptermodality{532}

\begin{chapteropening}
\textbf{By the end of this chapter:} \emph{I can talk about decorating and drilling at home.}

Complete the last lesson of chapter 532: I can talk about decorating and drilling at home.
\end{chapteropening}

\section[\texorpdfstring{\ru{обои} (obói)}{обои (obói)}]{\ru{обои} (obói) — wallpaper}

\label{lesson:RU-C532-oboi}



\begin{quote}
Before the new one: say the Russian for to fold, then the Russian for to throw away.
\end{quote}

\subsection*{You'll want to know: \ru{обои}}
\textbf{\ru{обои}} (\emph{obói}) — \textquotedblleft{}wallpaper\textquotedblright{}.

\begin{quote}
\ru{Мы} \ru{выбрали} \ru{светлые} \ru{обои} \ru{для} \ru{спальни}.
\end{quote}

\begin{quote}
\emph{My výbrali svétlyye obói dlya spálni.}
\end{quote}

\begin{quote}
We chose light wallpaper for the bedroom.
\end{quote}

Plural only; genitive \ru{обоев}.

\begin{cousinweb}
\emph{\ru{обои}} is usually traced to \emph{\ru{обить}} 'to cover, to upholster', from \emph{\ru{о}-} and \emph{\ru{бить}} 'to beat'; EtymDB lists \emph{\ru{обои}} without a source.
\end{cousinweb}

\begin{grammarlens}[title={What you've built}]
One more word for telling, explaining and describing at B1.
\end{grammarlens}

\subsection*{Your turn}
\begin{itemize}
\raggedright
  \item \emph{Say it:} \emph{obói}
  \item \emph{Say it:} \emph{obói}, once more
  \item \emph{Say it:} a sentence of your own with \emph{obói}
  \item \emph{Recall:} say the Russian for to fold, then the Russian for to throw away, then say \emph{obói} again
\end{itemize}

\begin{tcolorbox}[breakable,colback=teal!4,colframe=teal!35!black,title={Before you move on}]
What does \ru{обои} mean? (\textquotedblleft{}Wallpaper\textquotedblright{}.) Say it once more.
\end{tcolorbox}
\section[\texorpdfstring{\ru{клеить} (kléit)}{клеить (kléit)}]{\ru{клеить} (kléit) — to glue; to hang (wallpaper)}

\label{lesson:RU-C532-kleit}



\begin{quote}
Before the new one: say the Russian for to throw away, then the Russian for wallpaper.
\end{quote}

\subsection*{You'll want to know: \ru{клеить}}
\textbf{\ru{клеить}} (\emph{kléit}) — \textquotedblleft{}to glue; to hang (wallpaper)\textquotedblright{}.

\begin{quote}
\ru{В} \ru{выходные} \ru{мы} \ru{клеили} \ru{обои}.
\end{quote}

\begin{quote}
\emph{V vykhodnýye my kléili obói.}
\end{quote}

\begin{quote}
At the weekend we were hanging wallpaper.
\end{quote}

Imperfective; present \ru{клею}, \ru{клеишь}. Perfective \ru{склеить} or \ru{поклеить}.

\begin{cousinweb}
\emph{\ru{клеить}} is built on \emph{\ru{клей}} 'glue'; EtymDB lists the verb without a source.
\end{cousinweb}

\begin{grammarlens}[title={What you've built}]
One more word for telling, explaining and describing at B1.
\end{grammarlens}

\subsection*{Your turn}
\begin{itemize}
\raggedright
  \item \emph{Say it:} \emph{kléit}
  \item \emph{Say it:} \emph{kléit}, once more
  \item \emph{Say it:} a sentence of your own with \emph{kléit}
  \item \emph{Recall:} say the Russian for to throw away, then the Russian for wallpaper, then say \emph{kléit} again
\end{itemize}

\begin{tcolorbox}[breakable,colback=teal!4,colframe=teal!35!black,title={Before you move on}]
What does \ru{клеить} mean? (\textquotedblleft{}To glue; to hang (wallpaper)\textquotedblright{}.) Say it once more.
\end{tcolorbox}
\section[\texorpdfstring{\ru{сверлить} (sverlít)}{сверлить (sverlít)}]{\ru{сверлить} (sverlít) — to drill}

\label{lesson:RU-C532-sverlit}



\begin{quote}
Before the new one: say the Russian for wallpaper, then the Russian for to glue.
\end{quote}

\subsection*{You'll want to know: \ru{сверлить}}
\textbf{\ru{сверлить}} (\emph{sverlít}) — \textquotedblleft{}to drill\textquotedblright{}.

\begin{quote}
\ru{Сосед} \ru{с} \ru{утра} \ru{сверлит} \ru{стену}.
\end{quote}

\begin{quote}
\emph{Soséd s utrá sverlít sténu.}
\end{quote}

\begin{quote}
The neighbour has been drilling the wall since the morning.
\end{quote}

Imperfective; perfective \ru{просверлить}. The drill bit is a \emph{\ru{сверло}}.

\begin{cousinweb}
\emph{\ru{сверлить}} is built on \emph{\ru{сверло}} 'a drill bit', which is inherited from Proto-Slavic \emph{sv\ru{ь}rd\ru{ь}lo} (EtymDB); EtymDB has no entry for the verb.
\end{cousinweb}

\begin{grammarlens}[title={What you've built}]
One more word for telling, explaining and describing at B1.
\end{grammarlens}

\subsection*{Your turn}
\begin{itemize}
\raggedright
  \item \emph{Say it:} \emph{sverlít}
  \item \emph{Say it:} \emph{sverlít}, once more
  \item \emph{Say it:} a sentence of your own with \emph{sverlít}
  \item \emph{Recall:} say the Russian for wallpaper, then the Russian for to glue, then say \emph{sverlít} again
\end{itemize}

\begin{tcolorbox}[breakable,colback=teal!4,colframe=teal!35!black,title={Before you move on}]
What does \ru{сверлить} mean? (\textquotedblleft{}To drill\textquotedblright{}.) Say it once more.
\end{tcolorbox}
\section[\texorpdfstring{\ru{дрель} (drel)}{дрель (drel)}]{\ru{дрель} (drel) — a power drill}

\label{lesson:RU-C532-drel}



\begin{quote}
Before the new one: say the Russian for to glue, then the Russian for to drill.
\end{quote}

\subsection*{You'll want to know: \ru{дрель}}
\textbf{\ru{дрель}} (\emph{drel}) — \textquotedblleft{}a power drill\textquotedblright{}.

\begin{quote}
\ru{У} \ru{нас} \ru{есть} \ru{дрель}, \ru{но} \ru{нет} \ru{свёрл}.
\end{quote}

\begin{quote}
\emph{U nas yest drel, no net svyorl.}
\end{quote}

\begin{quote}
We've got a drill but no drill bits.
\end{quote}

Feminine.

\begin{cousinweb}
\emph{\ru{дрель}} is usually traced to Dutch or German, where \emph{drillen} means 'to bore'; EtymDB has no entry.
\end{cousinweb}

\begin{grammarlens}[title={What you've built}]
One more word for telling, explaining and describing at B1.
\end{grammarlens}

\subsection*{Your turn}
\begin{itemize}
\raggedright
  \item \emph{Say it:} \emph{drel}
  \item \emph{Say it:} \emph{drel}, once more
  \item \emph{Say it:} a sentence of your own with \emph{drel}
  \item \emph{Recall:} say the Russian for to glue, then the Russian for to drill, then say \emph{drel} again
\end{itemize}

\begin{tcolorbox}[breakable,colback=teal!4,colframe=teal!35!black,title={Before you move on}]
What does \ru{дрель} mean? (\textquotedblleft{}A power drill\textquotedblright{}.) Say it once more.
\end{tcolorbox}
\section[\texorpdfstring{\ru{отвёртка} (otvyórtka)}{отвёртка (otvyórtka)}]{\ru{отвёртка} (otvyórtka) — a screwdriver}

\label{lesson:RU-C532-otvyortka}



\begin{quote}
Before the new one: say the Russian for to drill, then the Russian for a power drill.
\end{quote}

\subsection*{You'll want to know: \ru{отвёртка}}
\textbf{\ru{отвёртка}} (\emph{otvyórtka}) — \textquotedblleft{}a screwdriver\textquotedblright{}.

\begin{quote}
\ru{Папа} \ru{ищет} \ru{отвёртку} \ru{в} \ru{ящике}.
\end{quote}

\begin{quote}
\emph{Pápa íshchet otvyórtku v yáshchike.}
\end{quote}

\begin{quote}
Dad is looking for the screwdriver in the drawer.
\end{quote}

Feminine.

\begin{cousinweb}
\emph{\ru{отвёртка}} is built on \emph{\ru{отвертеть}}, \emph{\ru{отвернуть}} 'to unscrew', from \emph{\ru{вертеть}} 'to turn'; EtymDB has no entry.
\end{cousinweb}

\begin{grammarlens}[title={What you've built}]
One more word for telling, explaining and describing at B1.
\end{grammarlens}

\subsection*{Your turn}
\begin{itemize}
\raggedright
  \item \emph{Say it:} \emph{otvyórtka}
  \item \emph{Say it:} \emph{otvyórtka}, once more
  \item \emph{Say it:} a sentence of your own with \emph{otvyórtka}
  \item \emph{Recall:} say the Russian for to drill, then the Russian for a power drill, then say \emph{otvyórtka} again
\end{itemize}

\begin{tcolorbox}[breakable,colback=teal!4,colframe=teal!35!black,title={Before you move on}]
What does \ru{отвёртка} mean? (\textquotedblleft{}A screwdriver\textquotedblright{}.) Say it once more.
\end{tcolorbox}
\section[(dialogue)]{Words from chapters 525–532 — a review}

\label{lesson:RU-R532-review-525-532}



\begin{quote}
Say the Russian for a poet and a poem.
\end{quote}

\subsection*{Your turn}
Say these aloud:

\begin{itemize}
\raggedright
  \item a poet, a poem, a rhyme, to rhyme, by heart
  \item a cover, a bookmark, to leaf through, a publishing house, to reread
  \item a TV series, an episode, subtitles, to switch, a screen adaptation
  \item a games console, a video game, a blogger, a subscriber, a gamer
  \item dust, cleaning, to vacuum, a cloth, to wipe
  \item laundry, a clothes peg, a stain, to hang out, to wash out
  \item a mess, junk, a storeroom, to fold, to throw away
  \item wallpaper, to glue, to drill, a power drill, a screwdriver
\end{itemize}

\begin{tcolorbox}[breakable,colback=teal!4,colframe=teal!35!black,title={Before you move on}]
A poet? (\textbf{\ru{поэт}}.) A poem? (\textbf{\ru{стихотворение}}.) A rhyme? (\textbf{\ru{рифма}}.) To rhyme? (\textbf{\ru{рифмовать}}.) By heart? (\textbf{\ru{наизусть}}.) A cover? (\textbf{\ru{обложка}}.) A bookmark? (\textbf{\ru{закладка}}.) To leaf through? (\textbf{\ru{листать}}.) A publishing house? (\textbf{\ru{издательство}}.) To reread? (\textbf{\ru{перечитывать}}.) A TV series? (\textbf{\ru{сериал}}.) An episode? (\textbf{\ru{эпизод}}.) Subtitles? (\textbf{\ru{субтитры}}.) To switch? (\textbf{\ru{переключать}}.) A screen adaptation? (\textbf{\ru{экранизация}}.) A games console? (\textbf{\ru{приставка}}.) A video game? (\textbf{\ru{видеоигра}}.) A blogger? (\textbf{\ru{блогер}}.) A subscriber? (\textbf{\ru{подписчик}}.) A gamer? (\textbf{\ru{геймер}}.) Dust? (\textbf{\ru{пыль}}.) Cleaning? (\textbf{\ru{уборка}}.) To vacuum? (\textbf{\ru{пылесосить}}.) A cloth? (\textbf{\ru{тряпка}}.) To wipe? (\textbf{\ru{вытирать}}.) Laundry? (\textbf{\ru{стирка}}.) A clothes peg? (\textbf{\ru{прищепка}}.) A stain? (\textbf{\ru{пятно}}.) To hang out? (\textbf{\ru{развешивать}}.) To wash out? (\textbf{\ru{отстирывать}}.) A mess? (\textbf{\ru{беспорядок}}.) Junk? (\textbf{\ru{хлам}}.) A storeroom? (\textbf{\ru{кладовка}}.) To fold? (\textbf{\ru{складывать}}.) To throw away? (\textbf{\ru{выбрасывать}}.) Wallpaper? (\textbf{\ru{обои}}.) To glue? (\textbf{\ru{клеить}}.) To drill? (\textbf{\ru{сверлить}}.) A power drill? (\textbf{\ru{дрель}}.) A screwdriver? (\textbf{\ru{отвёртка}}.)
\end{tcolorbox}
